\chapter{Stress Case Studies}\label{ch:m1:stress-case-studies}

At 14:45:28 on 6 May 2010 the most liquid equity future in the world stopped
trading for five seconds. In the preceding four and a half minutes it had
fallen more than 5\%; resting buy orders, above a hundred thousand contracts
that morning, were down to about a thousand. When trading resumed the price
turned, and within twenty minutes most of the fall was gone, except in three
hundred shares and funds where, meanwhile, twenty thousand trades had printed
at prices more than 60\% from where they had been: a cent, or a hundred
thousand dollars. Nothing in the economy changed between 14:32 and 15:00.
Every chapter of this book describes a mechanism that works on an ordinary
day. This one looks at four days on which mechanisms met their limits, from
the official reports, and asks each time the same three questions: what flow
arrived, who was supposed to absorb it, and why did they step back?

\section{6 May 2010}

\begin{definition}[Flash crash and stub quote]\label{def:m1:stress-case-studies:flash}
A \emph{flash crash}\index{flash crash} is a fall and recovery of prices, of a
size normally seen over days, within minutes, without corresponding news. A
\emph{stub quote}\index{stub quote} is a quote at a price far from the market
(a bid of a cent, an offer of \$100\,000), entered by or for a \omterm{def:m1:what-a-trading-firm-does:market-maker}{market maker} to
satisfy an obligation to quote continuously while not intending to trade.
\end{definition}

The joint report of the two US regulators sets out the sequence. The market
was already down and nervous. At 14:32 a mutual fund complex began selling
75\,000 E-mini contracts, about \$4.1 billion, as a hedge, through an
algorithm told to sell 9\% of the previous minute's volume, ``without regard
to price or time''. A comparable programme by the same trader, using
algorithms that took price and time into account, had earlier taken more
than five hours for its first 75\,000 contracts; this one finished in twenty
minutes.

High-frequency firms and other intermediaries bought the first contracts,
then, their inventory limits reached, sold them on. Between 14:45:13 and
14:45:27 they traded over 27\,000 contracts, 49\% of the volume, while buying
about 200 net: the report's ``hot-potato'' effect. The volume this generated
told the algorithm to sell faster. By 14:45:28 buy-side depth was under 1\,050
contracts, less than 1\% of its level at the start of the day. The exchange's
stop-logic functionality then paused trading for five seconds, and in that
pause buyers returned.

In the share market the same withdrawal had a different end. \omterm{def:m1:what-a-trading-firm-does:market-maker}{Market makers}
who pulled their quotes left only their \omterm{def:m1:stress-case-studies:flash}{stub quotes} behind; market orders
and stop-loss orders converted to market orders executed against them. The
exchanges later cancelled the trades that were more than 60\% away from the
14:40 prices, a threshold chosen after the fact, which was the last of the
day's uncertainties for anyone who had bought on the way down.

\begin{omfigure}
\begin{tikzpicture}
\begin{axis}[width=11cm,height=5.2cm,
  xlabel={minutes since the programme starts},ylabel={price change (\%)},
  tick label style={font=\small},label style={font=\small},
  xmin=1,xmax=40,ymin=-11,ymax=0.5,grid=major,grid style={gray!25},
  legend columns=3,legend style={font=\footnotesize,at={(0.5,-0.3)},anchor=north,draw=none,fill=none,/tikz/every even column/.append style={column sep=8pt}},
  table/col sep=comma]
\addplot[omExo,very thick,dotted] table[x=minute,y=none]{figdata/markets-1/31-stress-case-studies/paths.csv};
\addplot[omThm,very thick] table[x=minute,y=feedback]{figdata/markets-1/31-stress-case-studies/paths.csv};
\addplot[omDef,very thick,dashed] table[x=minute,y=pause]{figdata/markets-1/31-stress-case-studies/paths.csv};
\legend{a fixed book,depth withdraws and volume churns,the same with a pause at $-5\%$}
\end{axis}
\end{tikzpicture}

\omcaption{A toy model of the mechanism, calibrated to nothing: a programme
sells 9\% of the last minute's volume; depth shrinks, and intermediaries'
churn raises volume, when the price is below its recent average. With a fixed
book the same 75\,000 contracts cost 2.2\%; with the two feedbacks 9.8\%, of
which 5.9 points in the last four minutes; a pause that lets buyers re-anchor
stops the fall at 7.0\%. The model has no recovery. Data: the tutorial's simulation.}\label{fig:m1:stress-case-studies:paths}
\end{omfigure}

\section{24 August 2015}

\begin{definition}[Limit up--limit down and market-wide circuit breaker]\label{def:m1:stress-case-studies:luld}
\emph{Limit up--limit down}\index{limit up--limit down} is the US mechanism,
introduced after 2010, that prevents trades in a security outside a band
around a moving reference price and pauses trading if the market stays at the
band. A \emph{market-wide circuit breaker}\index{market-wide circuit breaker}
halts all equity trading when a broad index falls by set percentages in a day
(\cref{ch:m1:exchange-traded-funds}).
\end{definition}

\begin{dated}{2026-09}{The band rules, as the regulator's staff describe them}\label{dat:m1:stress-case-studies:luld}
The reference price is the mean price of trades over the preceding five
minutes; a new one takes effect only if it differs by 1\% or more from the
last, and stays at least 30 seconds. Bands are 5\%, 10\% or 20\% either side,
by tier of security and price, doubled from 09:30 to 09:45 and from 15:35 to
16:00, and multiplied by the leverage ratio for leveraged products. When the
best offer reaches the lower band, or the best bid the upper, the security is
in a limit state; if that lasts 15 seconds the primary exchange declares a
pause of at least five minutes, on all venues.
\end{dated}

On Monday 24 August 2015, after falls in Asia, the E-mini reached its
overnight limit, 5\% down, and was paused from 09:25 to 09:30. The largest
index fund opened 5.2\% down, was 7.8\% down by 09:35, and closed 4.2\% down.
By 09:35 the largest listing exchange had opened only 38\% of its shares in
the main index, so that for some minutes the index was computed partly from
Friday's closing prices and the funds' arbitrageurs had
no basket to price (\cref{prop:m1:exchange-traded-funds:stale}). There were
1\,278 limit-up-limit-down halts that day: 1\,058 in 327 exchange-traded
products and 220 in 144 other securities. A fifth of exchange-traded
products fell 20\% or more at some point, against 5\% of ordinary shares,
although four fifths of them were never halted.

The mechanism designed after 2010 worked as specified and produced a new
failure. A fund whose \omterm{def:m1:what-a-trading-firm-does:market-maker}{market makers} cannot value its basket has no anchor; a
halt stops trading but does not supply the anchor; and a reopening auction in
an unanchored fund, with stop-loss orders queued, can reopen it at the band
and halt it again. Repeated halts in products whose underlying shares were
trading normally were the day's signature.

\begin{omfigure}
\begin{tikzpicture}
\begin{axis}[width=9.5cm,height=4.4cm,ybar,bar width=16pt,
  ylabel={count},
  symbolic x coords={exchange-traded products,other securities},xtick=data,
  tick label style={font=\small},label style={font=\small},xticklabel style={font=\small\strut},
  ymin=0,ymax=1300,ymajorgrids,grid style={gray!25},enlarge x limits=0.5,
  scaled y ticks=false,yticklabel style={/pgf/number format/1000 sep={\,}},
  nodes near coords,nodes near coords style={font=\footnotesize,fill=white,inner sep=1pt,/pgf/number format/1000 sep={\,}},
  legend columns=2,legend style={font=\footnotesize,at={(0.5,-0.32)},anchor=north,draw=none,fill=none,/tikz/every even column/.append style={column sep=8pt}},
  table/col sep=comma]
\addplot[fill=omThm!60,draw=omThm] table[x=group,y=halts]{figdata/markets-1/31-stress-case-studies/halts.csv};
\addplot[fill=omDef!60,draw=omDef] table[x=group,y=securities]{figdata/markets-1/31-stress-case-studies/halts.csv};
\legend{halts,securities halted}
\end{axis}
\end{tikzpicture}

\omcaption{Limit-up-limit-down halts on 24 August 2015. The funds that halted
did so three times each on average. Data: the regulator's research note of
December 2015.}\label{fig:m1:stress-case-studies:halts}
\end{omfigure}

\section{5 February 2018 and January 2021}

Two episodes already met. On 5 February 2018
(\cref{ex:m1:volatility-first-contact:feb2018}) products that promised the
inverse, or a multiple, of volatility futures had to buy those futures at the
close in proportion to the day's rise; their buying was the rise, and the
largest inverse product lost nearly all of its value in a session. The flow
was mechanical, public and concentrated in minutes: the rebalancing
arithmetic of \cref{prop:m1:exchange-traded-funds:rebalance} applied to a
market too small for it.

In January 2021 (\cref{ex:m1:stock-loan-and-short-selling:gme}) a share with a
\omterm{def:m1:stock-loan-and-short-selling:si}{short interest} of 123\% of its float rose 2\,700\% in three weeks. The
mechanisms under strain were not on any exchange. Option \omterm{def:m1:what-a-trading-firm-does:market-maker}{market makers} short
the calls bought by individuals hedged by buying shares
(\cref{def:m1:zero-day-options-and-retail-flow:dealer}); short sellers
covered; and the clearing house's margin model, seeing the volatility and the
one-sided positions of retail brokers' customers, raised those brokers'
requirements several-fold before dawn (\cref{ch:m1:clearing-and-settlement}, \cref{ch:m1:margin}). Some brokers could
not or would not fund the call and restricted their customers to closing
purchases. The staff report found that covering was a small part of the
buying and that the rise was sustained by sentiment; what the episode added
to the list of things that can stop a market is a retail broker's balance
sheet.

\begin{definition}[Liquidity spiral]\label{def:m1:stress-case-studies:spiral}
A \emph{liquidity spiral}\index{liquidity spiral} is a loop in which falling
prices reduce the capacity of intermediaries to hold positions (through
losses, \omterm{def:m1:financing:margincall}{margin calls}, risk limits or uncertainty about the validity of their
trades), their withdrawal reduces depth, and reduced depth makes the same
flow move prices further.
\end{definition}

\begin{definition}[Clearly erroneous trade]\label{def:m1:stress-case-studies:erroneous}
A \emph{clearly erroneous trade}\index{clearly erroneous trade} is one that an
exchange may cancel after the fact because its price was too far from a
reference price. On 6 May 2010 the threshold, 60\%, was chosen after the
event; price bands exist so that such trades do not happen in the first
place.
\end{definition}

\section{What the four have in common}

\begin{method}[Reading an incident]\label{met:m1:stress-case-studies:read}
\begin{enumerate}
\item \textbf{The flow.} Who had to trade, how much, by what rule, by when?
In each case the flow was insensitive to price: a volume-following algorithm,
stop-loss orders, a daily rebalance, a \omterm{def:m1:financing:margincall}{margin call}.
\item \textbf{The absorbers.} Who normally takes the other side, and what
limits them: inventory, capital, a hedge that must exist, the certainty that
trades will stand?
\item \textbf{The loop.} Through which variable does the price move feed back
into the flow or into the absorbers' capacity: volume, volatility, margin, a
stale valuation?
\item \textbf{The circuit.} What stopped it: a pause, a band, an auction, a
rule change, new capital? What did the stopping mechanism itself break?
\item \textbf{The aftermath.} Which trades stood? Who bore the loss? Which
rule changed?
\end{enumerate}
\end{method}

\begin{center}\small
\begin{tabular}{p{2.2cm}p{3.3cm}p{3.4cm}p{3.6cm}}
\toprule
 & Price-insensitive flow & Why absorbers withdrew & What changed afterwards \\
\midrule
6 May 2010 & volume-following sell algorithm & inventory limits; fear of broken trades & \omterm{def:m1:stress-case-studies:flash}{stub quotes} prohibited (November 2010); price bands \\
24 Aug 2015 & market and stop orders at the open & no basket price for funds; late openings & one harmonised reopening auction with widening collars (2017) \\
5 Feb 2018 & daily rebalance of volatility products & futures market too small for the flow & the largest inverse product redeemed by its issuer \\
Jan 2021 & retail call buying; covering & brokers' clearing margin & shorter \omterm{def:m1:clearing-and-settlement:cycle}{settlement cycle} (\cref{ch:m1:clearing-and-settlement}) \\
\bottomrule
\end{tabular}
\end{center}

For a trading firm the lessons are operational. Know which of your
strategies are absorbers and what makes each one stop; decide in advance
whether it stops by widening, by shrinking or by switching off, because in
the event there are four minutes. Know which exchange rules decide whether
your trades stand. Hold the cash for the \omterm{def:m1:financing:margincall}{margin call} that arrives with the
opportunity (\cref{pb:m1:margin:1}). And keep a replayable record
(\cref{ch:m1:reading-market-data}): a day that can be run again is a day that
can be learned from.

\begin{omfigure}
\begin{tikzpicture}
\begin{axis}[width=10cm,height=4.6cm,
  xlabel={largest fall tolerated (\%)},ylabel={depth needed (thousand)},
  tick label style={font=\small},label style={font=\small},
  xmin=1.5,xmax=8.5,ymin=0,ymax=190,grid=major,grid style={gray!25},
  nodes near coords,nodes near coords style={font=\footnotesize,above=2pt,fill=white,inner sep=1pt,/pgf/number format/.cd,fixed,precision=0},
  table/col sep=comma]
\addplot[omDef,very thick,mark=*] table[x=max_fall_pct,y=depth_k]{figdata/markets-1/31-stress-case-studies/depth_needed.csv};
\draw[thin,dashed] (axis cs:1.5,100) -- (axis cs:8.5,100);
\node[font=\footnotesize,anchor=north east] at (axis cs:8.4,98) {the depth there was};
\end{axis}
\end{tikzpicture}

\omcaption{In the toy model, the initial depth that would have kept the
programme's fall within a given size. With the feedbacks, 16\% more depth
turns a 9.8\% fall into 5\%, and 61\% more brings it to 2\%, close to the
fall with no feedback at all: stability is very non-linear in depth, which is
why it disappears so abruptly. Data: the
tutorial's simulation.}\label{fig:m1:stress-case-studies:depth}
\end{omfigure}

\section{Tutorial: a liquidity withdrawal}\label{tut:m1:stress-case-studies:withdrawal}

\begin{tutorial}
\textbf{Goal.} Reproduce the shape of a volume-following sell programme in a
book whose depth depends on the fall so far; add a pause; find the depth that
would have absorbed the flow. \textbf{End state:} the two simulated figures of
this chapter.
\begin{tutsteps}
\item \textbf{One minute of the model,} in sixty steps, so that the result
does not depend on the step. Sales from last minute's volume; depth and churn
from the fall below a five-minute average price; impact from sales over
depth.
\omcode{code/markets-1/31-stress-case-studies/python/withdrawal.py}{20}{36}{A sell programme, a withdrawing book and a pause.}\label{lst:m1:stress-case-studies:run}
\item \textbf{Switch the feedbacks off} one at a time, then slow the programme
to 3\% of volume, and see what does the damage.
\item \textbf{Bands and pauses} as a state machine over a tape.
\omcode{code/firm/replay/firm_replay.py}{47}{77}{Reference price, limit state and pause.}\label{lst:m1:stress-case-studies:luld}
\end{tutsteps}
\textbf{What to change next.} Give the programme a \omterm{def:m1:asian-and-emerging-equity-markets:limit}{price limit} (do not sell
more than 1\% below the arrival price) and compare its cost and duration.
Then add a recovery: buyers who return in proportion to the distance below a
slowly moving fair value.
\end{tutorial}

\section{Build: the incident replay harness}\label{bld:m1:stress-case-studies:replay}

\begin{build}
\bfield{Purpose} The miniature firm runs every strategy through a library of
bad days before it trades, and replays its own incidents afterwards. The
harness turns a tape into the timeline a post-mortem needs.
\bfield{Interface} \texttt{Tick(t, bid, ask, trade)}; \texttt{Luld(band\_pct,
\ldots)} with \texttt{pct(t)}, \texttt{bands(t)}, \texttt{on\_tick(tick)} and a
list of \texttt{events}; \texttt{replay(ticks, luld)} returning the state and
bands at each tick; \texttt{depth\_within(levels, mid, bp)}.
\bfield{Rules} The band rules of \cref{dat:m1:stress-case-studies:luld},
parameterised: window, minimum move, minimum hold, doubling periods, limit
duration, pause length. After a pause the first trade sets a new reference.
Inputs come from the feed normaliser of
\cref{bld:m1:reading-market-data:feed}, so that a recorded day can be
replayed unmodified.
\bfield{Acceptance tests} \texttt{code/firm/replay/tests/}: doubling at the
open and the close; a reference that moves only by 1\% and after 30 seconds; a
limit state that becomes a pause after 15 seconds and a reopening; a limit
state that clears in time; depth near the mid.
\bfield{Stretch} A scenario library: for each of this chapter's four days, a
synthetic tape with the documented features, and a report of what each of the
firm's strategies would have done.
\end{build}

\begin{omsources}
\item US Commodity Futures Trading Commission and US Securities and Exchange
Commission, \emph{Findings Regarding the Market Events of May 6, 2010}, 30
September 2010.
\item US Securities and Exchange Commission, Division of Trading and Markets,
\emph{Research Note: Equity Market Volatility on August 24, 2015}, December
2015.
\item Bank for International Settlements, ``The equity market turbulence of 5
February: the role of exchange-traded volatility products'', \emph{BIS
Quarterly Review}, March 2018.
\item US Securities and Exchange Commission, \emph{Staff Report on Equity and
Options Market Structure Conditions in Early 2021}, 14 October 2021.
\item M.~Brunnermeier and L.~H.~Pedersen, ``Market liquidity and funding
liquidity'', \emph{Review of Financial Studies} 22 (2009).
\end{omsources}

\section{Exercises}

\begin{exercise}[$\star$]\label{exo:m1:stress-case-studies:1}
The E-mini programme of 6 May 2010 was 75\,000 contracts worth \$4.1 billion.
Give the implied index level (multiplier \$50). Between 14:32 and 14:45 it sold
about 35\,000 contracts: what share of the programme, and at what average
rate per minute?
\end{exercise}

\begin{exercise}[$\star$]\label{exo:m1:stress-case-studies:2}
A share's reference price is 40.00 and its band 5\%. Give the bands at 11:00
and at 09:40. The best offer falls to 38.00 at 11:00:00 and stays there.
When is the pause declared, and until when does it last at least?
\end{exercise}

\begin{exercise}[$\star$]\label{exo:m1:stress-case-studies:3}
From \cref{fig:m1:stress-case-studies:halts}: how many halts per halted
security, in exchange-traded products and in other securities? What does the
difference suggest?
\end{exercise}

\begin{exercise}[$\star\star$]\label{exo:m1:stress-case-studies:4}
An algorithm sells 9\% of the previous minute's volume. Volume is 20\,000
contracts a minute in calm conditions. How long does it take to sell 75\,000
contracts? If the algorithm's own selling and the churn it provokes triple
the volume, how long? Why is ``percentage of volume'' not a neutral
instruction?
\end{exercise}

\begin{exercise}[$\star\star$]\label{exo:m1:stress-case-studies:5}
A \omterm{def:m1:what-a-trading-firm-does:market-maker}{market maker} is obliged to quote two-sided at all times. Explain how a \omterm{def:m1:stress-case-studies:flash}{stub
quote} satisfies the letter of the obligation, what happened to market orders
that met \omterm{def:m1:stress-case-studies:flash}{stub quotes} on 6 May 2010, and what replaced \omterm{def:m1:stress-case-studies:flash}{stub quotes}.
\end{exercise}

\begin{exercise}[$\star\star$]\label{exo:m1:stress-case-studies:6}
On 24 August 2015 an index fund traded 20\% below the value of its holdings
while the holdings' own markets functioned. Using
\cref{prop:m1:exchange-traded-funds:band}, say which term of the arbitrage
band had become very large, and why a five-minute pause did not shrink it.
\end{exercise}

\begin{exercise}[$\star\star\star$]\label{exo:m1:stress-case-studies:7}
\emph{Coding.} With \texttt{run} and \texttt{trough} reproduce the three
falls of \cref{fig:m1:stress-case-studies:paths}. Then switch off only the
churn, and only the withdrawal, and run the full model at a participation of
3\%. Why does each feedback need the other to be dangerous?
\end{exercise}

\begin{exercise}[$\star\star\star$]\label{exo:m1:stress-case-studies:8}
\emph{Find the flaw.} A risk report: ``Our market-making strategy is safe in a
crash: it is flat at the end of every day, its worst day in five years of
backtest is $-0.4\%$ of capital, and it widens its quotes automatically when
volatility rises.'' Name three things the chapter's four days did that this
report cannot see.
\end{exercise}

\section{Problem: Five Seconds}

\begin{problem}[{Weekend problem --- how much depth would have been enough?}]\label{pb:m1:stress-case-studies:1}
Use the chapter's toy model with its default parameters: initial buy-side
depth 100\,000 contracts, calm volume 20\,000 a minute, a programme of 75\,000
contracts selling 9\% of the previous minute's volume, impact of 3\% for
selling a quantity equal to the current depth. With $f$ the fall of the price
below its average over the last five minutes or so, depth is multiplied by
$0.1 + 0.9\exp(-150 f)$ and calm volume by $1 + 100 f$.

\medskip\noindent\textbf{Part I --- The first minute.}
\begin{enumerate}
\item Give the programme's sales in the first minute.
\item Give the price change they cause if depth stays at its initial level.
\item When the price is 1\% below its recent average, give the depth and the
volume.
\item How small can depth become, and at what recent fall is it within a
tenth of that floor?
\item At a recent fall of 1\%, how much does the programme sell in a minute
and what does that minute cost?
\end{enumerate}

\medskip\noindent\textbf{Part II --- The whole path.}
\begin{enumerate}[resume]
\item The simulation gives falls of 2.2\% with a fixed book and 9.8\% with the
feedbacks. How many times more costly is the second?
\item Value the difference on 75\,000 contracts at \$54\,667 each, taking the
average sale at half the trough.
\item With a pause at $-5\%$ after which buyers measure the fall from the
reopening price, the trough is 7.0\%. What did the pause save?
\item In the model, why does the pause work? Relate it to what the official
report says happened in the five seconds.
\end{enumerate}

\medskip\noindent\textbf{Part III --- The depth that was missing.}
\begin{enumerate}[resume]
\item From \cref{fig:m1:stress-case-studies:depth}, give the initial depth
that keeps the fall within 2\%, and within 5\%.
\item Express each as a multiple of the actual 100\,000.
\item Why is the relation so non-linear?
\item The programme could instead have been given a participation rate of
3\%. Without running the model, what would change, and what would it cost the
seller?
\item Which party was best placed to prevent the event at the least cost?
\end{enumerate}

\medskip\noindent\textbf{Part IV --- Judgement.}
\begin{enumerate}[resume]
\item Intermediaries traded 27\,000 contracts in fourteen seconds and bought
200 net. Were they providing liquidity?
\item Should the 20\,000 trades at absurd prices have been cancelled? Give
the argument on each side.
\item Bands and pauses now exist. Name one thing they fixed and one they did
not, using 24 August 2015.
\item What should a volume-following execution algorithm always have?
\item State the \emph{named result}: the initial depth that would have
absorbed the programme with a fall of no more than 2\%.
\item In one sentence: what is liquidity, on the day it matters?
\end{enumerate}
\end{problem}

\section{Interview questions}

\begin{interviewq}[$\star$ \iqroles{trader, researcher, developer}]\label{iq:m1:stress-case-studies:1}
What happened on 6 May 2010?
\end{interviewq}

\begin{interviewq}[$\star$ \iqroles{trader, developer}]\label{iq:m1:stress-case-studies:2}
How does \omterm{def:m1:stress-case-studies:luld}{limit up--limit down} work?
\end{interviewq}

\begin{interviewq}[$\star\star$ \iqroles{trader, researcher}]\label{iq:m1:stress-case-studies:3}
Why did \omterm{def:m1:exchange-traded-funds:etf}{exchange-traded funds} trade far below the value of their holdings on
24 August 2015?
\end{interviewq}

\begin{interviewq}[$\star\star$ \iqroles{researcher, mle}]\label{iq:m1:stress-case-studies:4}
Your backtest covers five calm years. How do you assess the strategy's
behaviour in a crash?
\end{interviewq}

\begin{interviewq}[$\star\star$ \iqroles{trader, developer}]\label{iq:m1:stress-case-studies:5}
Your market-making system sees the price fall 3\% in two minutes on rising
volume. What should it do, and who decides?
\end{interviewq}

\begin{interviewq}[$\star\star\star$ \iqroles{researcher, trader, bank}]\label{iq:m1:stress-case-studies:6}
What do market crashes driven by market structure have in common, and what
would you look for to anticipate the next one?
\end{interviewq}
